% !PRINTED 425
\weberpage{425}

By the definition (11), $Q$ is an integral rational function
of $r$ which, when one orders it according to powers of $r$, has
integral rational numerical coefficients. If we therefore set
\[
r = r_0^k,
\]
then
\[
Q = F(r_0^k),
\]
where $F$ is the symbol for a rational function whose
coefficients do not depend on $k$. If we now set for
$k$ in turn the values
\[
k = 1,\ 2 \ldots n-1,
\]
then $r_0^k$ represents all the roots of the cyclotomic equation $X_n = 0$,
and the sum
\[
\overset{k}{\Sigma} Q = F(r_0) + F(r_0^2) + \cdots + F(r_0^{n-1})
\]
is a \gesperrt{symmetric function} of the roots of this equation; it
can therefore be expressed rationally in terms of the coefficients of the equation,
that is, in terms of rational numbers, and is consequently itself a
\gesperrt{rational number}. On the other hand, by the formula (15),
\[
\overset{k}{\Sigma} Q = h\, i^{\frac{n-1}{2}} \sqrt{n},
\]
where $h$ denotes an integer, namely the number of cases
in which the positive sign is to be taken in (15), diminished by
the number of cases in which the negative sign holds. But the two
are compatible with one another only when $h = 0$; and
with this the following theorem is proved:

\gesperrt{1. If $k$ runs through the sequence of numbers 1, 2, 3, \dots\ $n-1$,
then in the formula (16) the positive sign holds just as often as the
negative.}

It hardly needs special mention that one may also take for $k$
another complete system of residues of $n$, excluding the number
divisible by $n$.

\begin{center}
\S.\,136.\\[2pt]
\gesperrt{Primitive congruence roots.}
\end{center}

Parallel to the theory of the roots of unity runs a theory of the
so-called binomial congruences, without which, in the theory of the
cyclotomic equations, further steps cannot
%CONT>>
